% ---
% the following macOS command will strip all extended attributes
% like (Where from) from a file making the (user-specific metadata
% stand out more
% xattr -c /path/to/file.pdf
% ---

\documentclass[11pt]{article}

\usepackage[margin=0.75in]{geometry}
\usepackage[sfdefault]{roboto}
\usepackage[T1]{fontenc}
\usepackage{enumitem}
\usepackage{titlesec}
\usepackage{tabularx}
\usepackage[none]{hyphenat} % disable automatic hyphenation // ATS related
\usepackage[hidelinks]{hyperref} % hidden clickabble links

\pagestyle{empty}
\setlength{\parindent}{0pt}


% ---
% user-specific metadata
% ---
\usepackage[en-US]{datetime2} % sets default date style
\DTMsetdatestyle{iso}
\hypersetup{
    pdfauthor={Kevin Sikora},
    pdftitle={Kevin Sikora - Resume - \today},
    pdfsubject={Kevin Sikora's resume for a senior software engineer position.},
    pdfkeywords={Senior software engineer, Senior SDET, QA, Software testing, WordPress plugin development, Website development, Website maintenance, Drone pilot},
    pdfcreator={},
    pdfproducer={},
    pdfnewwindow=true
}


% ---
% company and position formatting
% ---
\newcounter{companycount}
\newcounter{positioncount}

\newcommand{\companyheader}[2]{%
    \ifnum\value{companycount}>0
        \par\addvspace{10pt}
    \fi
    \stepcounter{companycount}
    \setcounter{positioncount}{0}
    \noindent
    \begin{tabular*}{\textwidth}{@{}l@{\extracolsep{\fill}}r@{}}
        \textbf{#1} & #2 \\
    \end{tabular*}
    \par\nobreak
}

\newcommand{\positionheader}[2]{%
    \ifnum\value{positioncount}>0
        \par\addvspace{5pt}
    \fi
    \stepcounter{positioncount}
    \noindent
    \begin{tabular*}{\textwidth}{@{}l@{\extracolsep{\fill}}r@{}}
        \textit{#1} & #2 \\
    \end{tabular*}
    \par\nobreak
}


% ---
% optional custom commands
% ---
\newcommand{\websiteslabel}{\textit{Websites:}}
\newcommand{\githublabel}{\textit{GitHub:}}
\newcommand{\techlabel}{\textit{Technologies:}}
\newcommand{\italicizedlabel}[1]{\textit{#1:}} % \italicizedlabel{Some label}


% ---
% section heading formatting
% ---
\titleformat{\section}{\large\bfseries}{}{0em}{}[\titlerule]
\titlespacing{\section}{0pt}{22pt}{10pt}


% ---
% indent all bullet points
% ---
\setlist[itemize]{
    leftmargin=0.25in,
    itemsep=0pt,
    topsep=1pt
}


\begin{document} % specific resume information gets listed below


% ---
% Resume header
% ---
\begin{center}
    {\Huge \textbf{Your Name}} \\[6pt]
    {\small
    \href{mailto:your.email@example.com}{your.email@example.com}
    \, | \,
    (123) 456-7890
    \, | \,
    \href{https://www.linkedin.com/in/your-linkedin}{linkedin.com/in/your-linkedin}
    } \\[4pt]
    {\small Your Professional Title}
\end{center}


% ---
% Candidate summary
% ---
\section*{Summary}

Write a brief 2-4 sentence summary here. Focus on your main skills, years of experience, career focus, and the value you bring to a team.


% ---
% Work Experience
% ---
\section*{Work Experience}
\setcounter{companycount}{0} % used for spacing between the listed companies

% copy and paste this section for each company you worked at
\companyheader{Company Name}{City, State}
\positionheader{Job Title}{Month Year - Month Year}
\begin{itemize}
    \item Describe your responsibility, achievement, or impact using strong action words.
    \item Add numbers when possible, such as revenue, time saved, team size, or growth.
    \item \websiteslabel{} https://LiveTheDTOUR.com, https://CareerCurveball.com    
    \item \githublabel{} https://github.com/sikoracollective
    \item \techlabel{} Python, SQL, Tableau, Salesforce, Excel
\end{itemize}

\companyheader{Company Name}{City, State}
\positionheader{Job Title}{Month Year - Month Year}
\begin{itemize}
    \item Describe your work and the result it created.
    \item Highlight tools, processes, leadership, or business impact.
    \item \techlabel{} JavaScript, React, Node.js, Git, Jira
\end{itemize}

\companyheader{Company Name}{City, State}
\positionheader{Job Title 1}{Month Year - Month Year}
% copy this section to list different positions you had at the same company
\begin{itemize}
    \item Describe your work and the result it created.
    \item Highlight tools, processes, leadership, or business impact.
\end{itemize}
\positionheader{Job Title 2}{Month Year - Month Year}
\begin{itemize}
    \item Describe your work and the result it created.
    \item Highlight tools, processes, leadership, or business impact.
\end{itemize}


% ---
% Certifications
% ---
\section*{Certifications}

% copy and paste this section for each certification you have
\begin{itemize}
    \item \begin{tabularx}{\linewidth}{@{}X r@{}}
        \textbf{FAA Part 107 Remote Pilot Certificate} & February 2026 \\
    \end{tabularx}
    % copy and paste the following section as many times as need be
    \item \begin{tabularx}{\linewidth}{@{}X r@{}}
        \textbf{EASA UAS Remote Pilot Certificate} -- Drone Open Category A1+A3 & February 2026 \\
    \end{tabularx}
\end{itemize}


% ---
% Education
% ---
\section*{Education}

% copy and paste this section for each degree you have
\begin{tabular*}{\textwidth}{@{}l@{\extracolsep{\fill}}r@{}}
    \textbf{School Name} & City, State \\
    % example: \textbf{University of Wisconsin -- Madison} & Madison, Wisconsin \\
    % copy and paste this section as much as you need to if you received multiple
    % degrees from a university
    \textit{Degree 1 or Program Name 1} & Graduation Month Year \\
    % example: \textit{Bachelor of Science in Computer Science} &  June 2016 \\
    \textit{Degree 2 or Program Name 2} \\ % remove the month year if you don't want it
\end{tabular*}

\end{document}

